% Generated draft: requires mathematical and source review.
\documentclass[11pt]{article}
\usepackage[margin=25mm]{geometry}
\usepackage{fontspec}
\setmainfont{Latin Modern Roman}
\usepackage{amsmath,amssymb,mathtools,unicode-math}
\setmathfont{Latin Modern Math}
\usepackage{enumitem,needspace,xurl,hyperref}
\setlength{\parindent}{0pt}
\setlength{\parskip}{5pt}
\setlength{\emergencystretch}{3em}
\newcommand{\N}{\mathbb{N}}
\newcommand{\Z}{\mathbb{Z}}
\newcommand{\Q}{\mathbb{Q}}
\newcommand{\R}{\mathbb{R}}
\newcommand{\C}{\mathbb{C}}
\newcommand{\F}{\mathbb{F}}
\newcommand{\A}{\mathbb{A}}
\newcommand{\PP}{\mathbb{P}}
\newcommand{\E}{\mathbb{E}}
\newcommand{\eps}{\varepsilon}
\newcommand{\dd}{\,\mathrm{d}}
\begin{document}
\section*{UnsolvedMath ID: 30001043}
\textbf{Source problem code:} catalogue entry
\subsection{Definitions and mathematical statement}
Let $f:M\to M$ be the intended partially hyperbolic diffeomorphism with continuous $Df$-invariant splitting
\[
TM=E^s\oplus E^c\oplus E^u.
\]
In the standard compact-manifold convention, $E^s$ is uniformly contracted, $E^u$ uniformly expanded, and the splitting is dominated, with $E^s$ dominated by $E^c$ and $E^c$ by $E^u$. For example, after an iterate the relevant norm/conorm inequalities separate these three bundles uniformly; the extract does not specify its precise partial-hyperbolicity convention.

Retain the given hypothesis that $E^s$ and $E^u$ are uniquely integrable to $C^0$ foliations with smooth leaves. Unique integrability of $E^c$ means that through every point there is a unique local germ of an integral submanifold of dimension $\dim E^c$ tangent everywhere to $E^c$, with the resulting integral leaves forming the intended center foliation. A $C^0$ foliation describes transverse regularity and does not require its tangent distribution to be smooth transversely.

The question is whether $E^c$ is always uniquely integrable under exactly the stated general setup. No one-dimensional-center, bunching, accessibility, volume-preservation or symplectic hypothesis is added. Existence of some chosen invariant center foliation is not silently substituted for uniqueness of local integral manifolds; integrability of the center-stable and center-unstable sums is also a distinct condition.

The manifold class, differentiability, center dimension, exact local uniqueness and foliation regularity are not supplied. The literature assessment also mentions a ``corresponding Liouville-measure statement'', but no such assertion appears in either mathematical statement, so it is not turned into an additional conjecture or a hidden hypothesis. The dated open label is retained only as source curation.
\subsection{Short English statement}
Does the center bundle of every partially hyperbolic diffeomorphism integrate uniquely, given the stable and unstable foliations? The source does not specify extra dimension or regularity hypotheses, and its unrelated Liouville-measure phrase is not added to the target.
\subsection{Sources}
\begin{itemize}
\item Supplied report: \emph{Geometric Group Theory, Hyperbolic Dynamics and Symplectic Geometry} (2009), Oberwolfach Reports, DOI 10.4171/owr/2008/42, pp. 2337--2406; proposed year 2008. \url{https://doi.org/10.4171/owr/2008/42}.
\item Local UnsolvedMath registry, permanent ID 30001043, catalogue entry; created 25 August 2026, updated 16 September 2026. Its corrected\_verified note describes notation cleanup. Literature checked at \texttt{2026-08-21}: the assessment is supplied dated curation, not independently verified. It labels the general center-integrability question open and mentions a Liouville-measure assertion absent from the extracted mathematical statement. No additional external reference is supplied, and no current-status conclusion is inferred from that mismatch.
\end{itemize}
\end{document}
